\documentclass[11pt]{article}
\usepackage[a4paper,margin=28mm]{geometry}
\usepackage{amsmath,amssymb,amsthm}
\usepackage[colorlinks=true,linkcolor=blue,urlcolor=blue,citecolor=blue]{hyperref}
\newtheorem{theorem}{Theorem}
\newtheorem{lemma}{Lemma}
\newtheorem{corollary}[lemma]{Corollary}
\newtheorem{proposition}[lemma]{Proposition}
\newcommand{\Q}{\mathbb Q}
\newcommand{\Z}{\mathbb Z}
\newcommand{\F}{\mathbb F}
\newcommand{\R}{\mathbb R}
\newcommand{\C}{\mathbb C}
\DeclareMathOperator{\End}{End}
\DeclareMathOperator{\Gal}{Gal}
\DeclareMathOperator{\GL}{GL}
\DeclareMathOperator{\PGL}{PGL}
\DeclareMathOperator{\Lie}{Lie}
\title{Rational \(p\)-torsion of rank two on geometrically split
abelian surfaces that are simple over \(\Q\)}
\author{}
\date{21 October 2026}

\begin{document}
\maketitle

\begin{theorem}\label{thm:main}
Let \(A/\Q\) be an abelian surface that is simple over \(\Q\) and
geometrically split, and let \(p\geq5\) be a prime. Then
\[
 \dim_{\F_p}A(\Q)[p]\leq1.
\]
In particular, the groups \([5,5]\), \([5,10]\) and \([7,7]\) cannot
occur as the rational torsion subgroup of a geometrically split
genus-\(2\) Jacobian that is simple over \(\Q\).
\end{theorem}

Here ``geometrically split'' means that \(A_{\overline{\Q}}\) is
isogenous to a product of elliptic curves. We impose no restriction on
the degree of this isogeny. The theorem does not require a principal
polarization.

The hypothesis of simplicity over \(\Q\) is essential. A product of two
elliptic curves over \(\Q\), each with a rational point of order \(5\),
has \((\Z/5\Z)^2\) in its rational torsion, and each of the three groups
has an example split over \(\Q\) in the census. The bound \(p\geq5\) is
also needed: the census contains geometrically split Jacobians, simple
over \(\Q\), with torsion \([3,3]\) and \([3,9]\).
Section~\ref{sec:hypotheses} says where the proof uses each hypothesis.

\section{Notation and standard facts}\label{sec:notation}

Fix an embedding \(\overline{\Q}\subset\C\). For a number field
\(K\subset\overline{\Q}\), write \(G_K=\Gal(\overline{\Q}/K)\), and let
\(c\in G_{\Q}\) be complex conjugation. Let \(p\) be an odd prime, and
put
\[
 T=T_p(A),\qquad V=T\otimes_{\Z_p}\Q_p,\qquad A[p]=T/pT,
\]
with their actions of \(G_{\Q}\). Let \(\rho:G_{\Q}\to\GL(T)\) be the
action on \(T\), and \(\chi:G_{\Q}\to\F_p^\times\) the mod-\(p\)
cyclotomic character. Since \(\chi\) is surjective, it has order
\(p-1\); in particular \(\chi\neq\mathbf 1\). For a finite-dimensional
representation \(M\) of \(G_{\Q}\) over \(\F_p\), we write
\(M^{\mathrm{ss}}\) for its semisimplification and \(M^*\) for its
dual.

Let \(D=\End(A_{\overline{\Q}})\otimes\Q\). The group \(G_{\Q}\) acts
on \(D\) by \(\Q\)-algebra automorphisms,
\({}^\sigma\varphi=\sigma\circ\varphi\circ\sigma^{-1}\). An element of
\(D\) fixed by \(G_K\) is defined over \(K\). Let \(L\) be the smallest
field over which all geometric endomorphisms of \(A\) are defined. It
is a finite Galois extension of \(\Q\), and \(\Gamma=\Gal(L/\Q)\) acts
faithfully on \(D\). We use the following standard facts.

\begin{enumerate}
\item[(F1)] \emph{Tate modules.} The map
\(\End(A_{\overline{\Q}})\otimes\Z_p\to\End_{\Z_p}(T)\) is injective
\cite[Theorem I.10.15]{Milne}, applied over \(\overline{\Q}\). Hence
\(D\otimes_{\Q}\Q_p\to\End_{\Q_p}(V)\) is injective. It is
\(G_{\Q}\)-equivariant when \(G_{\Q}\) acts on \(\End_{\Q_p}(V)\) by
conjugation:
\(T_p({}^\sigma\varphi)=\rho(\sigma)\,T_p(\varphi)\,\rho(\sigma)^{-1}\).

\item[(F2)] \emph{Lattice independence.} Let \(W\) be a
finite-dimensional \(\Q_p\)-vector space with a continuous action of
\(G_{\Q}\). For a \(G_{\Q}\)-stable \(\Z_p\)-lattice \(\Lambda\subset W\),
the semisimplification of \(\Lambda/p\Lambda\) depends only on the
isomorphism class of \(W\). Indeed, the characteristic polynomials on
\(\Lambda/p\Lambda\) are the reductions of those on \(W\), and the
Brauer--Nesbitt theorem applies. In particular, a lattice may be
transported through an isomorphism of representations.

\item[(F3)] \emph{Polarization.} Every abelian variety over \(\Q\) has
a polarization over \(\Q\). Combined with the Weil pairing
\cite[Section I.13]{Milne}, it gives a nondegenerate alternating
\(G_{\Q}\)-equivariant pairing \(V\times V\to\Q_p(1)\), hence an
isomorphism
\[
 V\simeq V^*\otimes_{\Q_p}\Q_p(1).
\]
The polarization is used only on the rational Tate module, so its
degree may be divisible by \(p\).

\item[(F4)] \emph{Idempotents over \(\Q\).} If \(D\) contains an
idempotent \(e\neq0,1\) fixed by \(G_{\Q}\), then \(A\) is not simple
over \(\Q\). Indeed, \(e\) is defined over \(\Q\), so
\(ne\in\End(A)\) for some integer \(n\geq1\). Its image is an abelian
subvariety of \(A\) over \(\Q\). It is nonzero because \(e\neq0\). It
is proper because \(ne\circ n(1-e)=0\) and \(n(1-e)\neq0\), so \(ne\)
is not an isogeny.

\item[(F5)] \emph{Complex conjugation.} Since \(A\) is defined over
\(\R\), complex conjugation of coordinates gives an antiholomorphic
involution \(\mathcal C\) of \(A(\C)\) which is a group automorphism.
Write \(A(\C)=\Lie(A_{\C})/\Lambda\) with \(\Lambda=H_1(A(\C),\Z)\), so
that \(H=\Lambda\otimes\R\) is identified with \(\Lie(A_{\C})\), and let
\(J\) be the complex structure on \(H\). The map \(\mathcal C\) lifts
to an \(\R\)-linear involution \(c_H\) of \(H\), its differential at
\(0\), and
\[
 c_HJ=-Jc_H,
\]
because \(\mathcal C\) is antiholomorphic. Each \(\varphi\in D\) acts
on \(H\) by an \(\R\)-linear map \(\rho_H(\varphi)\) commuting with
\(J\): this is the analytic representation, equivalently the rational
representation on \(H_1(A(\C),\Q)\) tensored with \(\R\). The resulting
map \(\rho_H:D\otimes_{\Q}\R\to\End_{\R}(H)\) is injective. Finally, for \(\varphi\in\End(A_{\overline{\Q}})\) we have
\({}^c\varphi=\mathcal C\circ\varphi\circ\mathcal C\) on \(A(\C)\),
since both sides are morphisms (the right-hand side is \(\varphi\) with
complex-conjugated coefficients) that agree on \(A(\overline{\Q})\).
Therefore, by linearity,
\[
 \rho_H({}^c\varphi)=c_H\,\rho_H(\varphi)\,c_H^{-1}
 \qquad(\varphi\in D).
\]
\end{enumerate}

\section{Consequences of rational torsion of rank two}\label{sec:rank-two}

In this section \(p\) is an odd prime and
\(\dim_{\F_p}A(\Q)[p]\geq2\). Simplicity over \(\Q\) and geometric
splitting are not used here.

\begin{lemma}\label{lem:ss}
We have \(A[p]^{\mathrm{ss}}\simeq\mathbf 1^2\oplus\chi^2\). Every
composition factor of \(\End_{\F_p}(A[p])\), on which \(G_{\Q}\) acts
by conjugation, is one of \(\mathbf 1\), \(\chi\), \(\chi^{-1}\).
\end{lemma}

\begin{proof}
The subspace \(A(\Q)[p]=A[p]^{G_{\Q}}\) is a trivial submodule of
dimension at least \(2\). Thus \(\mathbf 1\) occurs in
\(A[p]^{\mathrm{ss}}\) with some multiplicity \(m\geq2\).

Transport the lattice
\(\operatorname{Hom}_{\Z_p}(T,\Z_p)\otimes_{\Z_p}\Z_p(1)\) into \(V\)
through the isomorphism of (F3). Its reduction is
\(A[p]^*\otimes\chi\). By (F2),
\[
 A[p]^{\mathrm{ss}}\simeq\bigl(A[p]^*\otimes\chi\bigr)^{\mathrm{ss}}
 =\bigl(A[p]^{\mathrm{ss}}\bigr)^*\otimes\chi.
\]
The multiplicity of \(\chi\) on the right equals the multiplicity of
\(\mathbf 1\) in \((A[p]^{\mathrm{ss}})^*\), which is \(m\). Hence
\(\chi\) occurs at least twice in \(A[p]^{\mathrm{ss}}\). Since
\(\chi\neq\mathbf 1\) and \(\dim A[p]=4\), it follows that \(m=2\) and
\(A[p]^{\mathrm{ss}}\simeq\mathbf 1^2\oplus\chi^2\).

For the second claim,
\(\End_{\F_p}(A[p])\simeq A[p]\otimes A[p]^*\). In the Grothendieck
group of finite-dimensional \(\F_p[G_{\Q}]\)-modules, the class of a
module is the sum of the classes of its composition factors. Tensor
product over \(\F_p\) is exact in each variable, so the class of
\(M\otimes N\) depends only on the classes of \(M\) and \(N\),
bilinearly. Since
\((A[p]^*)^{\mathrm{ss}}\simeq\mathbf 1^2\oplus\chi^{-2}\), we get
\[
 [\End_{\F_p}(A[p])]
 =\bigl(2[\mathbf 1]+2[\chi]\bigr)\bigl(2[\mathbf 1]+2[\chi^{-1}]\bigr)
 =8[\mathbf 1]+4[\chi]+4[\chi^{-1}].
\]
This determines the composition factors.
\end{proof}

This statement concerns composition factors. It does not assert that
\(A[p]\) or its endomorphism module is semisimple. Only
\(\chi\neq\mathbf 1\) was used, so Lemma~\ref{lem:ss} holds for every
odd \(p\).

\begin{lemma}\label{lem:saturated}
Let \(B\subseteq\End_{\Q_p}(V)\) be a \(G_{\Q}\)-stable
\(\Q_p\)-subspace, and put
\[
 B_{\Z_p}=B\cap\End_{\Z_p}(T).
\]
Then \(B_{\Z_p}\) is a \(G_{\Q}\)-stable \(\Z_p\)-lattice in \(B\), and
reduction gives a \(G_{\Q}\)-equivariant injection
\[
 B_{\Z_p}/pB_{\Z_p}\hookrightarrow\End_{\F_p}(A[p]).
\]
Consequently, for every \(G_{\Q}\)-stable lattice
\(\Lambda\subseteq B\), every composition factor of
\(\Lambda/p\Lambda\) is one of \(\mathbf 1\), \(\chi\),
\(\chi^{-1}\).
\end{lemma}

\begin{proof}
The module \(B_{\Z_p}\) is a submodule of the finite free
\(\Z_p\)-module \(\End_{\Z_p}(T)\), hence finite free. It spans \(B\),
because \(p^kX\in\End_{\Z_p}(T)\) for every \(X\in B\) and \(k\) large.
It is \(G_{\Q}\)-stable because \(B\) and \(\End_{\Z_p}(T)\) are. It is
saturated in \(\End_{\Z_p}(T)\): if \(X\in B_{\Z_p}\) and
\(X\in p\End_{\Z_p}(T)\), then
\(p^{-1}X\in B\cap\End_{\Z_p}(T)=B_{\Z_p}\). Thus
\(B_{\Z_p}\cap p\End_{\Z_p}(T)=pB_{\Z_p}\), which is the injectivity
of the reduction map.

By the Jordan--H\"older theorem, the composition factors of
\(B_{\Z_p}/pB_{\Z_p}\) are among those of \(\End_{\F_p}(A[p])\), which
are listed in Lemma~\ref{lem:ss}. By (F2), \(\Lambda/p\Lambda\) and
\(B_{\Z_p}/pB_{\Z_p}\) have the same semisimplification.
\end{proof}

\begin{corollary}\label{cor:quadratic}
Suppose that \(p\geq5\) and \(\dim_{\F_p}A(\Q)[p]\geq2\). Let
\(\varepsilon:G_{\Q}\to\{\pm1\}\) be a nontrivial character. Then there
is no nonzero \(\delta\in D\) with
\({}^\sigma\delta=\varepsilon(\sigma)\delta\) for all
\(\sigma\in G_{\Q}\).
\end{corollary}

\begin{proof}
Suppose such a \(\delta\) exists. By (F1), its image in
\(\End_{\Q_p}(V)\) is nonzero, and \(B=\Q_p\delta\) is a
\(G_{\Q}\)-stable line on which \(G_{\Q}\) acts through
\(\varepsilon\). The lattice \(\Z_p\delta\) reduces to the character
\[
 \overline\varepsilon:G_{\Q}\to\{\pm1\}\subset\F_p^\times,
\]
which is nontrivial because \(-1\neq1\) in \(\F_p\). By
Lemma~\ref{lem:saturated}, \(\overline\varepsilon\) is one of
\(\mathbf 1\), \(\chi\), \(\chi^{-1}\). It is not \(\mathbf 1\). It has
order \(2\), whereas \(\chi\) and \(\chi^{-1}\) have order
\(p-1\geq4\). This is a contradiction.
\end{proof}

From now on, \(A\) is simple over \(\Q\) and geometrically split,
\(p\geq5\), and we suppose, to reach a contradiction, that
\(\dim_{\F_p}A(\Q)[p]\geq2\). Exactly one of the following holds:
\(A_{\overline{\Q}}\) is isogenous to \(E_1\times E_2\) with
\(E_1,E_2\) not isogenous; \(A_{\overline{\Q}}\) is isogenous to
\(E^2\) with \(E\) a CM elliptic curve; \(A_{\overline{\Q}}\) is
isogenous to \(E^2\) with \(E\) without CM. We treat these cases in
Sections~\ref{sec:nonisogenous}, \ref{sec:cm-square}, and
\ref{sec:endomorphism-field}--\ref{sec:complex-conjugation}.

\section{Nonisogenous geometric factors}\label{sec:nonisogenous}

Suppose that \(A_{\overline{\Q}}\) is isogenous to \(E_1\times E_2\),
where \(E_1\) and \(E_2\) are not geometrically isogenous. Since
\(\operatorname{Hom}(E_1,E_2)=0\),
\[
 D\simeq\End^0(E_1)\times\End^0(E_2),
\]
where each factor is \(\Q\) or an imaginary quadratic field. The only
idempotents of this product of two fields other than \(0\) and \(1\)
are \(e_1=(1,0)\) and \(e_2=(0,1)\). The group \(G_{\Q}\) acts by
algebra automorphisms, so it permutes \(\{e_1,e_2\}\). If it fixed
\(e_1\), then \(A\) would not be simple over \(\Q\), by (F4). Hence the
permutation action defines a nontrivial character
\(\varepsilon:G_{\Q}\to\{\pm1\}\), with \(\varepsilon(\sigma)=-1\)
exactly when \(\sigma\) exchanges \(e_1\) and \(e_2\). The element
\[
 \delta=e_1-e_2
\]
satisfies \({}^\sigma\delta=\varepsilon(\sigma)\delta\). This
contradicts Corollary~\ref{cor:quadratic}.

This argument also covers elliptic factors with complex
multiplication; only the two idempotents are used. Simplicity over
\(\Q\) is used to show that \(\varepsilon\) is nontrivial.

\section{The square of a CM elliptic curve}\label{sec:cm-square}

Suppose that \(A_{\overline{\Q}}\) is isogenous to \(E^2\), where \(E\)
has complex multiplication by an imaginary quadratic field \(M\). Then
\(D\simeq M_2(M)\), and its centre \(Z=Z(D)\) is an imaginary quadratic
field isomorphic to \(M\). Choose \(\delta\in Z\) with
\(\delta^2=-d\) for a rational number \(d>0\). Since \(G_{\Q}\) acts on
\(D\) by algebra automorphisms, it preserves \(Z\) and acts on it by
field automorphisms. Hence there is a character
\(\varepsilon:G_{\Q}\to\{\pm1\}\) with
\[
 {}^\sigma\delta=\varepsilon(\sigma)\delta\qquad(\sigma\in G_{\Q}).
\]

We show that \(\varepsilon(c)=-1\). Choose an isogeny
\(f:A_{\C}\to E_{\C}^2\). Conjugation by \(f\) identifies \(D\) with
\(\End^0(E^2)=M_2(\End^0(E))\). The field \(\End^0(E)\) acts on the
one-dimensional space \(\Lie(E_{\C})\) through a field embedding into
\(\C\), so the centre of \(M_2(\End^0(E))\) acts on
\(\Lie(E_{\C}^2)=\Lie(E_{\C})^2\) by complex scalars. Transporting
through the \(\C\)-linear isomorphism
\(df:\Lie(A_{\C})\to\Lie(E^2_{\C})\), we find a field embedding
\(\iota:Z\to\C\) such that
\[
 \rho_H(z)=\iota(z)\qquad(z\in Z),
\]
where the right-hand side means multiplication by the complex number
\(\iota(z)\) on \(H=\Lie(A_{\C})\). By (F5), and since \(c_H\) is
\(\C\)-antilinear,
\[
 \rho_H({}^c\delta)=c_H\,\iota(\delta)\,c_H^{-1}
 =\overline{\iota(\delta)}=-\iota(\delta)=\rho_H(-\delta).
\]
Here \(\overline{\iota(\delta)}=-\iota(\delta)\) because
\(\iota(\delta)^2=-d<0\). Since \(\rho_H\) is injective,
\({}^c\delta=-\delta\). Thus \(\varepsilon\) is nontrivial, and
Corollary~\ref{cor:quadratic} gives a contradiction.

Simplicity over \(\Q\) is not used in this case.

\section{The square of a non-CM elliptic curve: the endomorphism field}
\label{sec:endomorphism-field}

Suppose finally that \(A_{\overline{\Q}}\) is isogenous to \(E^2\),
where \(E\) has no complex multiplication. Then
\[
 D\simeq M_2(\Q).
\]
By Skolem--Noether, every \(\Q\)-algebra automorphism of \(M_2(\Q)\) is
inner. The faithful action of \(\Gamma\) on \(D\) therefore identifies
\(\Gamma\) with a finite subgroup of
\(\operatorname{Aut}(M_2(\Q))=\PGL_2(\Q)\).

\begin{lemma}\label{lem:pgl2}
A finite subgroup of \(\PGL_2(\Q)\) is isomorphic to one of
\[
 C_1,\ C_2,\ C_3,\ C_4,\ C_6,\ D_4,\ D_6,\ D_8,\ D_{12},
\]
where the subscript on a dihedral group denotes its order. In
particular, \(D_4=C_2\times C_2\), and the order of such a group divides
\(24\).
\end{lemma}

\begin{proof}
This is a special case of \cite[Introduction and
Proposition 1.1]{Beauville}. For completeness: a finite subgroup of
\(\PGL_2(\R)\) acts by isometries on the hyperbolic plane, so it fixes
a point. Its action on the tangent plane at that point embeds it in
\(O(2)\), so it is cyclic or dihedral. If an element of \(\PGL_2(\Q)\)
has order \(n>2\), the ratio of the eigenvalues of a representative
\(M\) is a primitive \(n\)-th root of unity \(\zeta\). Thus
\[
 \frac{(\operatorname{tr}M)^2}{\det M}=2+\zeta+\zeta^{-1}\in\Q,
\]
which forces \(\varphi(n)=2\), that is, \(n\in\{3,4,6\}\).
\end{proof}

\begin{proposition}\label{prop:cyclotomic}
In the present case,
\[
 (p,\Gamma)\in\{(5,C_4),\ (7,C_6)\},\qquad L=\Q(\zeta_p).
\]
\end{proposition}

\begin{proof}
Since \(p\geq5\) and \(|\Gamma|\) divides \(24\), the prime \(p\) does
not divide \(|\Gamma|\).

Let \(B\subseteq\End_{\Q_p}(V)\) be the image of
\(D\otimes_{\Q}\Q_p\), which is isomorphic to it by (F1), and let
\(B_{\Z_p}\) be the lattice of Lemma~\ref{lem:saturated}. The group
\(G_{\Q}\) acts on \(B\) through \(\Gamma\), and this action is
faithful because the action on \(D\) is. Put
\(\overline B=B_{\Z_p}/pB_{\Z_p}\).

\emph{The action of \(\Gamma\) on \(\overline B\) is faithful.} Let
\(g\in\Gamma\) act trivially on \(\overline B\), and suppose
\(g\neq1\). Choosing a basis of \(B_{\Z_p}\), write \(g=1+p^kY\) with
\(k\geq1\) and \(Y\) a matrix over \(\Z_p\) that is nonzero modulo
\(p\). Let \(m\) be the order of \(g\); it divides \(|\Gamma|\), so
\(p\nmid m\). Since \(2k\geq k+1\),
\[
 1=g^m\equiv1+mp^kY\pmod{p^{k+1}},
\]
which is impossible because \(m\) is a unit in \(\Z_p\). Hence
\(g=1\).

\emph{\(\Gamma\) is abelian, and \(\overline B\) is a sum of
characters.} By Maschke's theorem, \(\overline B\) is a semisimple
\(\F_p[\Gamma]\)-module, so it is the direct sum of its composition
factors. By Lemma~\ref{lem:saturated}, each of these is one of
\(\mathbf 1\), \(\chi\), \(\chi^{-1}\). Thus
\[
 \overline B\simeq\psi_1\oplus\psi_2\oplus\psi_3\oplus\psi_4,\qquad
 \psi_i\in\{\mathbf 1,\chi,\chi^{-1}\},
\]
where each \(\psi_i\) is trivial on \(G_L\). Since \(\Gamma\) acts
faithfully, it embeds in \((\F_p^\times)^4\), so \(\Gamma\) is abelian.

\emph{The cyclotomic field.} If \(\Gamma\) were trivial, every element
of \(D\) would be defined over \(\Q\), including the idempotent
\(\left(\begin{smallmatrix}1&0\\0&0\end{smallmatrix}\right)\). This
contradicts simplicity over \(\Q\), by (F4). Hence \(\Gamma\neq1\), and
by faithfulness some \(\psi_i\) is nontrivial, so it is \(\chi\) or
\(\chi^{-1}\). This character is trivial on \(G_L\), so
\(G_L\subseteq\ker\chi=G_{\Q(\zeta_p)}\), that is,
\(\Q(\zeta_p)\subseteq L\). Hence \(\Gamma\) has a cyclic quotient
\(\Gal(\Q(\zeta_p)/\Q)\) of order \(p-1\geq4\).

\emph{The list.} The abelian groups in Lemma~\ref{lem:pgl2} are
\(C_1,C_2,C_3,C_4,C_6\) and \(D_4=C_2\times C_2\). Their cyclic
quotients have orders dividing \(1,2,3,4,6\) and \(2\), respectively.
A cyclic quotient of even order \(p-1\geq4\) occurs only for
\(\Gamma=C_4\) with \(p=5\) and for \(\Gamma=C_6\) with \(p=7\). In both
cases \([L:\Q]=|\Gamma|=p-1=[\Q(\zeta_p):\Q]\), so \(L=\Q(\zeta_p)\).
\end{proof}

\section{Complex conjugation and cyclic endomorphism fields}
\label{sec:complex-conjugation}

Let \(\operatorname{sgn}:\PGL_2(\R)\to\{\pm1\}\) be the sign of the
determinant of a representative. It is well defined, since
\(\det(\lambda M)=\lambda^2\det M\), and it is a homomorphism. Every
isomorphism \(D\otimes_{\Q}\R\simeq M_2(\R)\) gives an isomorphism
\(\operatorname{Aut}(D\otimes_{\Q}\R)\simeq\PGL_2(\R)\), unique up to
conjugation, so the sign of an element of
\(\Gamma\subseteq\operatorname{Aut}(D)\subseteq
\operatorname{Aut}(D\otimes_{\Q}\R)\) is well defined.

\begin{lemma}\label{lem:sign}
An element of finite order of \(\PGL_2(\R)\) with negative determinant
has order \(2\). Consequently, if \(\Gamma\) is cyclic of order greater
than \(2\), every element of \(\Gamma\) has sign \(+1\).
\end{lemma}

\begin{proof}
Scale a representative so that \(\det M=-1\). Its characteristic
polynomial \(X^2-tX-1\) has discriminant \(t^2+4>0\), so \(M\) has real
eigenvalues \(\lambda\neq\mu\) with \(\lambda\mu=-1\). If \(M^n\) is
scalar, then \(\lambda^n=\mu^n\). Since \(\lambda\) and \(\mu\) are
real of opposite signs, this gives \(\mu=-\lambda\), hence
\(\lambda^2=1\) and \(M^2=I\). The class of \(M\) is not trivial
because \(\det M<0\), so it has order \(2\).

If \(\Gamma=\langle g\rangle\) has order greater than \(2\), then
\(\operatorname{sgn}(g)=+1\) by the first part, and so every power of
\(g\) has sign \(+1\).
\end{proof}

\begin{lemma}\label{lem:real}
In the situation of Section~\ref{sec:endomorphism-field}, the class in
\(\PGL_2(\R)\) of the action of \(c\) on \(D\otimes_{\Q}\R\) is either
trivial or has sign \(-1\). Consequently, if \(\Gamma\) is cyclic of
order greater than \(2\), then \(L\) is totally real.
\end{lemma}

\begin{proof}
Let \(H\), \(J\), \(c_H\) and \(\rho_H\) be as in (F5). The image
\(\rho_H(D\otimes_{\Q}\R)\) is a copy of \(M_2(\R)\) inside
\(\End_{\R}(H)\simeq M_4(\R)\). Every module over \(M_2(\R)\) is a sum
of copies of the standard module \(U=\R^2\). Hence there is a
two-dimensional real vector space \(W\) with
\[
 H\simeq U\otimes_{\R}W,\qquad
 \rho_H(D\otimes_{\Q}\R)=\End_{\R}(U)\otimes1.
\]
The commutant of \(\End_{\R}(U)\otimes1\) in \(\End_{\R}(U\otimes W)\)
is \(1\otimes\End_{\R}(W)\). Since \(J\) commutes with \(\rho_H(D)\),
we have \(J=1\otimes J_W\) with \(J_W^2=-1\).

By (F5), conjugation by \(c_H\) preserves \(\End_{\R}(U)\otimes1\) and
induces on it the action of \(c\) on \(D\otimes_{\Q}\R\). By
Skolem--Noether, this automorphism is conjugation by \(M\otimes1\) for
some \(M\in\GL(U)\), and the class \([M]\in\PGL_2(\R)\) is the image of
\(c\). Then \((M^{-1}\otimes1)c_H\) commutes with
\(\End_{\R}(U)\otimes1\), so
\[
 c_H=M\otimes N\qquad\text{for some }N\in\GL(W).
\]
The relation \(c_HJ=-Jc_H\) gives \(M\otimes(NJ_W+J_WN)=0\), so
\(NJ_W=-J_WN\). In a basis of \(W\) with
\(J_W=\left(\begin{smallmatrix}0&-1\\1&0\end{smallmatrix}\right)\),
this means
\[
 N=\begin{pmatrix}a&b\\b&-a\end{pmatrix},\qquad N^2=(a^2+b^2)I,
\]
with \(a^2+b^2>0\) because \(N\) is invertible. Since \(c_H^2=1\), we
get \((a^2+b^2)M^2\otimes1=1\otimes1\), so \(M^2=sI\) with
\(s=(a^2+b^2)^{-1}>0\). If \([M]\) is not trivial, then \(M\) is not
scalar, its minimal polynomial is \(X^2-s\), and its eigenvalues are
\(\sqrt s\) and \(-\sqrt s\). Thus \(\det M=-s<0\).

Now let \(\Gamma\) be cyclic of order greater than \(2\). By
Lemma~\ref{lem:sign}, the image of \(c\) in \(\Gamma\) has sign \(+1\),
so it is trivial by the first part. Since \(\Gamma\) acts faithfully on
\(D\), the restriction of \(c\) to \(L\) is trivial, so
\(L\subset\R\). As \(L/\Q\) is Galois, \(L\) is totally real.
\end{proof}

\paragraph{Completion of the proof.}
In Section~\ref{sec:endomorphism-field}, \(\Gamma\) is cyclic of order
\(p-1\geq4\) and \(L=\Q(\zeta_p)\). By Lemma~\ref{lem:real}, \(L\) is
totally real. But \(\Q(\zeta_p)\) is not real, a contradiction.
Together with Sections~\ref{sec:nonisogenous} and~\ref{sec:cm-square},
this shows that \(\dim_{\F_p}A(\Q)[p]\geq2\) is impossible, and proves
Theorem~\ref{thm:main}.
\hfill\(\square\)

\section{Where the hypotheses are used}\label{sec:hypotheses}

\paragraph{The bound \(p\geq5\).} Lemmas~\ref{lem:ss}
and~\ref{lem:saturated} hold for every odd prime. The bound \(p\geq5\)
is used in exactly three places.
\begin{enumerate}
\item In Corollary~\ref{cor:quadratic}, which serves
Sections~\ref{sec:nonisogenous} and~\ref{sec:cm-square}: a nontrivial
quadratic character is neither \(\chi\) nor \(\chi^{-1}\), because
\(\chi\) has order \(p-1\geq4\).
\item In Proposition~\ref{prop:cyclotomic}: \(p\nmid|\Gamma|\), so
that Maschke's theorem applies and \(\Gamma\) acts faithfully on
\(\overline B\).
\item In Proposition~\ref{prop:cyclotomic}: the cyclic quotient of
order \(p-1\geq4\) leaves only \(\Gamma=C_4\) and \(\Gamma=C_6\), which
are cyclic of order greater than \(2\).
\end{enumerate}

The theorem is false for \(p=3\). The census contains the geometrically
split Jacobians
\href{https://genus-2-torsion.github.io/curve.html?id=3.3.c}{3.3.c},
with torsion
\href{https://genus-2-torsion.github.io/group.html?g=3.3}{\([3,3]\)},
and \href{https://genus-2-torsion.github.io/curve.html?id=3.9.c}{3.9.c},
with torsion
\href{https://genus-2-torsion.github.io/group.html?g=3.9}{\([3,9]\)},
both certified simple over \(\Q\). For \(p=3\), the character
\(\chi_3\) is itself the quadratic character of \(\Q(\sqrt{-3})\), so
the first use of \(p\geq5\) listed above breaks down when \(\varepsilon\) is this character. This is what
happens for 3.9.c. The extended genus-\(2\) data of
Booker--Sutherland~\cite{BookerSutherland} list the geometric
endomorphism algebra \(\Q\times\Q\) for it (label 6804.a.1), and its
census splitting certificate is defined over \(\Q(\sqrt{-3})\). So the
character \(\varepsilon\) of Section~\ref{sec:nonisogenous} is
\(\chi_3\).

For 3.3.c, the second use listed above breaks down. With \(Y=2y+x^3\), its model becomes
\(Y^2=x^6-52x^3+784\), which has the geometric automorphisms
\[
 (x,Y)\mapsto(\zeta_3x,\,Y),\qquad
 (x,Y)\mapsto\bigl(\sqrt[3]{784}/x,\ 28Y/x^3\bigr).
\]
The first is defined over \(\Q(\zeta_3)\) and the second over
\(\Q(\sqrt[3]{784})\), and over no smaller fields. Automorphisms of a
curve of genus \(2\) act faithfully on its Jacobian, so \(L\) contains
\(\Q(\zeta_3,\sqrt[3]{784})\). This field has Galois group \(S_3\) over
\(\Q\), so \(3\) divides \(|\Gamma|\). The same data list
\(D=M_2(\Q)\) and Sato--Tate group \(J(E_3)\) for 3.3.c (label
1764.d.2), so it falls under Section~\ref{sec:endomorphism-field}.

\paragraph{Simplicity over \(\Q\).} It is used only twice: in
Section~\ref{sec:nonisogenous}, to make \(\varepsilon\) nontrivial, and
in Proposition~\ref{prop:cyclotomic}, to exclude \(\Gamma=1\). The
proof therefore shows, without assuming simplicity over \(\Q\), that
\(\dim_{\F_p}A(\Q)[p]\leq1\) for \(p\geq5\) whenever
\(A_{\overline{\Q}}\) is isogenous to the square of a CM elliptic
curve, or to the square of a non-CM elliptic curve with
\(\Gamma\neq1\). The case \(\Gamma=1\) genuinely occurs: the census
curve \href{https://genus-2-torsion.github.io/curve.html?id=5.5.a}{5.5.a},
with torsion \([5,5]\), is isogenous over \(\Q\) to the product of the
elliptic curves 11a1 and 11a3, which are isogenous over \(\Q\) and have
no CM. So \(D=M_2(\Q)\) and all its elements are defined over \(\Q\).

\paragraph{What is not used.} The proof uses no bound on the degree of
a geometric splitting isogeny and no assumption on the integral
endomorphism ring. The polarization may have degree divisible by
\(p\). A principal polarization is not assumed. The Weil pairing alone
does not exclude rational \(p\)-torsion of rank two, as the examples
split over \(\Q\) show. The proof uses the full Galois action on
\(A[p]\) and on the geometric endomorphism algebra.

\section{Consequence for the census}

A group \(A(\Q)_{\mathrm{tors}}\) of one of the forms \([5,5]\),
\([5,10]\) or \([7,7]\) has \(\dim_{\F_p}A(\Q)[p]=2\) for \(p=5\) or
\(p=7\). Hence the groups
\href{https://genus-2-torsion.github.io/group.html?g=5.5}{\([5,5]\)},
\href{https://genus-2-torsion.github.io/group.html?g=5.10}{\([5,10]\)}
and
\href{https://genus-2-torsion.github.io/group.html?g=7.7}{\([7,7]\)}
are impossible in the class ``geometrically split, simple over \(\Q\)''.
No other census group contains \((\Z/p\Z)^2\) for a prime \(p\geq5\).

Simplicity over \(\Q\) is essential. Each of the three groups has a
certified example split over \(\Q\):
\href{https://genus-2-torsion.github.io/curve.html?id=5.5.a}{5.5.a},
\href{https://genus-2-torsion.github.io/curve.html?id=5.10.a}{5.10.a}
and
\href{https://genus-2-torsion.github.io/curve.html?id=7.7.a}{7.7.a}.
The exclusion therefore concerns the ``geometrically split, simple over
\(\Q\)'' column only. It does not change the combined geometrically
split \(\infty\)? grade of any of the three groups, and none of them
becomes \(\emptyset\) on the geometrically split side. The theorem says
nothing about geometrically simple Jacobians; for example, \([5,5]\) has
the geometrically simple example
\href{https://genus-2-torsion.github.io/curve.html?id=5.5.b}{5.5.b}.

The theorem applies to every geometric splitting degree, including
degrees divisible by \(p\). It does not require a principal
polarization.

\begin{thebibliography}{9}
\bibitem{Beauville}
A. Beauville,
\emph{Finite subgroups of \(\PGL_2(K)\)}, in \emph{Vector bundles and
complex geometry}, Contemp. Math. \textbf{522}, Amer. Math. Soc., 2010,
23--29, Introduction and Proposition 1.1.
\href{https://arxiv.org/abs/0909.3942}{arXiv:0909.3942}.

\bibitem{Milne}
J. S. Milne,
\emph{Abelian varieties}, course notes, version 2.00 (2008), Chapter I,
Theorem 10.15 and Section 13.
\href{https://www.jmilne.org/math/CourseNotes/AV.pdf}{www.jmilne.org/math/CourseNotes/AV.pdf}.

\bibitem{BookerSutherland}
A. R. Booker and A. V. Sutherland,
\emph{Genus 2 curves of small conductor}, in preparation; preliminary
data in the alpha version of the LMFDB (curves 1764.d.2 and 6804.a.1).
\href{https://alpha.lmfdb.org/Genus2Curve/Q/}{alpha.lmfdb.org/Genus2Curve/Q/}.
Used only for the remark on \(p=3\) in Section~\ref{sec:hypotheses}.
\end{thebibliography}

\end{document}
